\section{E1: same-runtime supplementary comparison}
\label{app:e1}
E1 compares frozen H+guard and B-HPA on the previously observed
\path{locked_test_steady_v3} and \path{locked_test_recovery_v3}
inputs without revising the V3/V4 locks. All 110 formal lifecycle generations
share the two-A100 Qwen2.5-7B-Instruct FP16 runtime, model closure, service lookup,
health probe, synthetic API, and accounting contract. Both instances start
physically ready with empty queues; adapter targets remain B-HPA=0 and H+guard=2.
These targets and the intrinsic dispatch rules belong to the complete policy
treatments, so E1 does not isolate the causal contribution of the guard module.

\subsection{Identities, feasibility, and retained smoke outcomes}
The formal matrix contains one all2 capacity run per window and three fresh
repeats of each comparison policy per window: 2 capacity plus 12 comparison
identities. Each steady run contains 1,444 online and 120 offline arrivals;
each recovery run contains 334 online and 120 offline arrivals. Independent raw
recomputation finds all 14,126 formal arrivals completed on time, with zero late,
failed, or unfinished requests. Both capacity runs pass $G_F$ and every comparison
run passes P1/P1+. Table~\ref{tab:e1runs} retains all 14 identities; the capacity
runs are qualification evidence, not extra B-HPA or guard repeats.

\begin{table*}[t]
\centering
\small
\caption{All E1 formal identities. Occupied time is window GPU seconds; costs
are scenario USD. API is the percentage of all online arrivals, p99 is each
run's completed-online end-to-end quantile in seconds, and S/D counts window
starts/shutdowns. Every row completes all online and offline arrivals on time.}
\label{tab:e1runs}
\setlength{\tabcolsep}{3pt}
\par\vspace{4pt}
\begin{tabular*}{0.92\linewidth}{@{\extracolsep{\fill}}lllrrrrrr@{}}
\toprule
Window & Policy & Repeat & Occupied & Window cost & Full cost & API (\%) & p99 & S/D\\
\midrule
% Generated from independently reviewed P5_RAW_v2; see P6_TABLE_MANIFEST.json.
steady & all2 & gate & 4200.0 & 4.2000 & 4.2700 & 0.00 & 1.616 & 0/0\\
steady & B-HPA & 1 & 2407.0 & 2.4678 & 2.5389 & 6.37 & 19.002 & 2/3\\
steady & B-HPA & 2 & 2407.1 & 2.4696 & 2.5398 & 6.58 & 19.005 & 2/3\\
steady & B-HPA & 3 & 2407.2 & 2.4680 & 2.5383 & 6.37 & 19.003 & 2/3\\
steady & H+guard & 1 & 1266.0 & 1.7611 & 1.8339 & 53.25 & 19.002 & 13/15\\
steady & H+guard & 2 & 1140.5 & 1.6153 & 1.6842 & 50.83 & 19.002 & 13/15\\
steady & H+guard & 3 & 1140.3 & 1.6151 & 1.6794 & 50.83 & 19.002 & 13/15\\
recovery & all2 & gate & 4200.0 & 4.2000 & 4.2758 & 0.00 & 1.586 & 0/0\\
recovery & B-HPA & 1 & 2044.6 & 2.0715 & 2.1415 & 11.68 & 15.002 & 1/2\\
recovery & B-HPA & 2 & 2044.4 & 2.0713 & 2.1386 & 11.68 & 15.002 & 1/2\\
recovery & B-HPA & 3 & 2044.4 & 2.0713 & 2.1439 & 11.68 & 15.002 & 1/2\\
recovery & H+guard & 1 & 1102.2 & 1.2576 & 1.3279 & 60.18 & 15.008 & 12/14\\
recovery & H+guard & 2 & 794.5 & 0.9447 & 1.0110 & 58.98 & 15.008 & 11/13\\
recovery & H+guard & 3 & 794.4 & 0.9446 & 1.0104 & 58.98 & 15.008 & 11/13\\
\bottomrule

\end{tabular*}
\end{table*}

The 17 entries include three 90-second nonformal smokes.
B-HPA a1 is \texttt{TECHNICAL\_INVALID}: its lifecycle callback rejected
ready after forced stopping, leaving four offline arrivals unfinished.
After repair, a2 records 68/68 online completions on time but 0/4 offline
completions at cutoff; its \texttt{NEGATIVE\_RESULT} and false gate remain intact.
The wrapper's \texttt{PASS} denotes technical integration only.
Guard's smoke completes all 68 online and four offline arrivals on time.
None enters the formal comparison. The raw \texttt{formal\_accepted=false}
field is not overwritten; scientific qualification is independently derived from
request and lifecycle records.

\subsection{Resource use and cost boundary}
For window $w$, let $\overline C_{g,w}$ and $\overline C_{h,w}$ be the arithmetic
means across the three H+guard and B-HPA repeats. E1 reports
$R_w=1-\overline C_{g,w}/\overline C_{h,w}$ only after both arms qualify in
all three runs, then $R=(R_{\rm steady}+R_{\rm recovery})/2$.
This is the equal average of within-window ratios, not a pooled ratio or the
average of six per-pair percentages. The fixed repeat blocks support descriptive
paired differences, not confidence intervals or population-level significance.
Window scenario-cost reductions are 32.5975\% and 49.3598\%, hence 40.9786\%
equally weighted. Full-deployment reductions are 31.7650\% and 47.8629\%.

\begin{table}[t]
\centering
\small
\caption{E1 mean phase costs (scenario USD), three repeats per comparison
arm/window. Full includes setup, window, and cleanup.}
\label{tab:e1phases}
\setlength{\tabcolsep}{2pt}
\par\vspace{4pt}
\begin{tabular*}{\linewidth}{@{\extracolsep{\fill}}llrrrr@{}}
\toprule
Window & Policy & Setup & Window & Cleanup & Full\\
\midrule
% Generated from independently reviewed P5_RAW_v2; see P6_TABLE_MANIFEST.json.
steady & B-HPA & 0.0675 & 2.4685 & 0.0031 & 2.5390\\
steady & H+guard & 0.0687 & 1.6638 & 0.0000 & 1.7325\\
recovery & B-HPA & 0.0670 & 2.0714 & 0.0030 & 2.1413\\
recovery & H+guard & 0.0675 & 1.0489 & 0.0000 & 1.1164\\
\bottomrule

\end{tabular*}
\end{table}

Occupied time spans \texttt{launch\_issued} through \texttt{verified\_release},
including starting and stopping, clipped to each phase. Window occupied-time
reductions are 50.8839\% and 56.1240\%; the corresponding mean savings are
1,224.82 and 1,147.44 GPU seconds. This is local resource occupancy, not total
compute across the local and external services. Active time is not GPU compute
utilization. Guard's zero cleanup cost means that release already completed
within the window, not that release or setup is free.

Repricing the fixed ledgers makes the cost boundary explicit. Let
$\Delta C_G$, $\Delta C_A$, and $\Delta C_E$ denote mean H+guard minus
B-HPA GPU, API, and lifecycle-event costs in scenario USD.
Their steady values are $(-1.2248,0.3862,0.0340)$;
recovery gives $(-1.1474,0.0930,0.0320)$.
With GPU price multiplier $x$, a common API input/output-budget multiplier
$y$, and fixed lifecycle-event prices, the cost difference is
$\Delta C(x,y)=x\Delta C_G+y\Delta C_A+\Delta C_E$.
Neither arm incurs missed-job penalties. Since $\Delta C_A>0$, H+guard
costs less when
\[
y<y^*(x)=-\frac{x\Delta C_G+\Delta C_E}{\Delta C_A}.
\]
At $x=1$, the API break-even multipliers are 3.08 and 11.99
(3.09 and 12.02 for full deployment).
At $y=1$, window savings require GPU price multipliers above 0.343 and
0.109. Lower local GPU occupancy therefore need not imply lower cost:
the boundary depends on external-service prices. This analysis holds
routing and lifecycle events fixed, without retuning either policy.

\subsection{Routing, lifecycle, and queue trade-offs}
The online API fraction uses all online arrivals as denominator, not only local
completions. Its mean rises from 6.44\% to 51.64\% in steady and from 11.68\% to
59.38\% in recovery. Table~\ref{tab:e1tradeoffs} retains the additional lifecycle
events and queue exposure. Queue exposure sums each request's actual
release-to-dispatch time, or release-to-cutoff time if undispatched, in request
seconds. It is neither a queue-latency p99 nor a new value of the historical
prediction loss $L$. Tail latency uses completed online requests; the all-arrival
feasibility checks above independently rule out hidden formal failures.

\begin{table}[t]
\centering
\small
\caption{E1 trade-offs. S and D are mean window starts and shutdowns; Q is
mean queue request-seconds. p99 ranges span three per-run completed-online
quantiles (seconds), not confidence intervals or pooled-request quantiles.}
\label{tab:e1tradeoffs}
\setlength{\tabcolsep}{2pt}
\par\vspace{4pt}
\begin{tabular*}{\linewidth}{@{\extracolsep{\fill}}llrrrr@{}}
\toprule
Window & Policy & S & D & Q & p99 range\\
\midrule
% Generated from independently reviewed P5_RAW_v2; see P6_TABLE_MANIFEST.json.
steady & B-HPA & 2.00 & 3.00 & 1674.9 & 19.002--19.005\\
steady & H+guard & 13.00 & 15.00 & 12884.7 & 19.002--19.002\\
recovery & B-HPA & 1.00 & 2.00 & 751.1 & 15.002--15.002\\
recovery & H+guard & 11.33 & 13.33 & 9429.1 & 15.008--15.008\\
\bottomrule

\end{tabular*}
\end{table}

The two known inputs and three repeats per cell establish a bounded policy
comparison. Greater API use and queue exposure rule out all-dimensional
dominance; quality-equivalent API substitution, energy, provider bills, and
new-trace generalization remain unmeasured.

\subsection{Independent recomputation and artifact links}
The accepted table source is
\path{docs/maxopt_v4_r3_p5_20260929/ANALYSIS/P5_RAW_v2/}:
\path{FORMAL_RUNS.csv} retains 14 formal identities, \path{ATTEMPTS.csv}
all attempts, \path{PAIRS.csv} six repeat blocks, and \path{AGGREGATES.json}
the qualified estimands. \path{COMPARABILITY.json} binds plans, runtime,
and permitted policy differences. In the parent documentation directory,
\path{REVIEW_B.json} and \path{REVIEW_B_V2.json} bind accepted output hashes;
\path{P6_GENERATE_E1_TABLES.py} and \path{P6_TABLE_MANIFEST.json} reproduce
and bind the displayed rows. The standalone E1 package reproduces the resource
and price-boundary calculations from portable reduced records; its
\path{verify.py} checks the manifest and full-precision reference results.
